% Mode Counter (CTR): counter dienkripsi lalu di-XOR dengan plainteks.
% Dirujuk dari section-02.tex dengan % Mode Counter (CTR): counter dienkripsi lalu di-XOR dengan plainteks.
% Dirujuk dari section-02.tex dengan % Mode Counter (CTR): counter dienkripsi lalu di-XOR dengan plainteks.
% Dirujuk dari section-02.tex dengan % Mode Counter (CTR): counter dienkripsi lalu di-XOR dengan plainteks.
% Dirujuk dari section-02.tex dengan \input{figures/mode-ctr}.
\begin{figure}[htbp]
    \centering
    \begin{tikzpicture}[
        kotak/.style={draw, rounded corners, align=center, font=\small, inner sep=3pt,
                      minimum width=1.35cm, minimum height=0.6cm},
        blok/.style={kotak, fill=blue!6},
        enc/.style={kotak, fill=orange!12, minimum width=1.1cm},
        cip/.style={kotak, fill=green!10},
        xor/.style={draw, circle, inner sep=1pt, font=\small, fill=yellow!18,
                    minimum size=0.5cm},
        grup/.style={draw, dashed, gray!60, rounded corners, inner sep=5pt},
        panah/.style={-{Latex[length=2mm]}, thick},
    ]
        \foreach \i/\x/\t in {1/1.9/{ctr}, 2/4.9/{ctr$+1$}, 3/7.9/{ctr$+2$}} {
            \node[blok] (k\i) at (\x,0)    {\t};
            \node[enc]  (e\i) at (\x,-1.2) {$E_K$};
            \node[xor]  (x\i) at (\x,-2.4) {$\oplus$};
            \node[cip]  (c\i) at (\x,-3.7) {$C_{\i}$};
            \node[blok] (p\i) at (\x-2.0,-2.4) {$P_{\i}$};
            \draw[panah] (k\i) -- (e\i);
            \draw[panah] (e\i) -- (x\i);
            \draw[panah] (p\i) -- (x\i);
            \draw[panah] (x\i) -- (c\i);
            % Kotak putus-putus menandai satu blok yang diproses mandiri,
            % supaya jelas P_i hanya menuju XOR pada panelnya sendiri.
            \node[grup, fit=(k\i)(c\i)(p\i)] {};
        }
        \node[font=\scriptsize, align=center] at (4.9,-4.5)
            {tidak ada ketergantungan antar blok: seluruh blok dapat diproses secara paralel};
    \end{tikzpicture}
    \caption{Mode Counter (CTR): nilai counter yang terus bertambah dienkripsi lalu di-XOR dengan plainteks}
    \label{fig:mode-ctr}
\end{figure}
.
\begin{figure}[htbp]
    \centering
    \begin{tikzpicture}[
        kotak/.style={draw, rounded corners, align=center, font=\small, inner sep=3pt,
                      minimum width=1.35cm, minimum height=0.6cm},
        blok/.style={kotak, fill=blue!6},
        enc/.style={kotak, fill=orange!12, minimum width=1.1cm},
        cip/.style={kotak, fill=green!10},
        xor/.style={draw, circle, inner sep=1pt, font=\small, fill=yellow!18,
                    minimum size=0.5cm},
        grup/.style={draw, dashed, gray!60, rounded corners, inner sep=5pt},
        panah/.style={-{Latex[length=2mm]}, thick},
    ]
        \foreach \i/\x/\t in {1/1.9/{ctr}, 2/4.9/{ctr$+1$}, 3/7.9/{ctr$+2$}} {
            \node[blok] (k\i) at (\x,0)    {\t};
            \node[enc]  (e\i) at (\x,-1.2) {$E_K$};
            \node[xor]  (x\i) at (\x,-2.4) {$\oplus$};
            \node[cip]  (c\i) at (\x,-3.7) {$C_{\i}$};
            \node[blok] (p\i) at (\x-2.0,-2.4) {$P_{\i}$};
            \draw[panah] (k\i) -- (e\i);
            \draw[panah] (e\i) -- (x\i);
            \draw[panah] (p\i) -- (x\i);
            \draw[panah] (x\i) -- (c\i);
            % Kotak putus-putus menandai satu blok yang diproses mandiri,
            % supaya jelas P_i hanya menuju XOR pada panelnya sendiri.
            \node[grup, fit=(k\i)(c\i)(p\i)] {};
        }
        \node[font=\scriptsize, align=center] at (4.9,-4.5)
            {tidak ada ketergantungan antar blok: seluruh blok dapat diproses secara paralel};
    \end{tikzpicture}
    \caption{Mode Counter (CTR): nilai counter yang terus bertambah dienkripsi lalu di-XOR dengan plainteks}
    \label{fig:mode-ctr}
\end{figure}
.
\begin{figure}[htbp]
    \centering
    \begin{tikzpicture}[
        kotak/.style={draw, rounded corners, align=center, font=\small, inner sep=3pt,
                      minimum width=1.35cm, minimum height=0.6cm},
        blok/.style={kotak, fill=blue!6},
        enc/.style={kotak, fill=orange!12, minimum width=1.1cm},
        cip/.style={kotak, fill=green!10},
        xor/.style={draw, circle, inner sep=1pt, font=\small, fill=yellow!18,
                    minimum size=0.5cm},
        grup/.style={draw, dashed, gray!60, rounded corners, inner sep=5pt},
        panah/.style={-{Latex[length=2mm]}, thick},
    ]
        \foreach \i/\x/\t in {1/1.9/{ctr}, 2/4.9/{ctr$+1$}, 3/7.9/{ctr$+2$}} {
            \node[blok] (k\i) at (\x,0)    {\t};
            \node[enc]  (e\i) at (\x,-1.2) {$E_K$};
            \node[xor]  (x\i) at (\x,-2.4) {$\oplus$};
            \node[cip]  (c\i) at (\x,-3.7) {$C_{\i}$};
            \node[blok] (p\i) at (\x-2.0,-2.4) {$P_{\i}$};
            \draw[panah] (k\i) -- (e\i);
            \draw[panah] (e\i) -- (x\i);
            \draw[panah] (p\i) -- (x\i);
            \draw[panah] (x\i) -- (c\i);
            % Kotak putus-putus menandai satu blok yang diproses mandiri,
            % supaya jelas P_i hanya menuju XOR pada panelnya sendiri.
            \node[grup, fit=(k\i)(c\i)(p\i)] {};
        }
        \node[font=\scriptsize, align=center] at (4.9,-4.5)
            {tidak ada ketergantungan antar blok: seluruh blok dapat diproses secara paralel};
    \end{tikzpicture}
    \caption{Mode Counter (CTR): nilai counter yang terus bertambah dienkripsi lalu di-XOR dengan plainteks}
    \label{fig:mode-ctr}
\end{figure}
.
\begin{figure}[htbp]
    \centering
    \begin{tikzpicture}[
        kotak/.style={draw, rounded corners, align=center, font=\small, inner sep=3pt,
                      minimum width=1.35cm, minimum height=0.6cm},
        blok/.style={kotak, fill=blue!6},
        enc/.style={kotak, fill=orange!12, minimum width=1.1cm},
        cip/.style={kotak, fill=green!10},
        xor/.style={draw, circle, inner sep=1pt, font=\small, fill=yellow!18,
                    minimum size=0.5cm},
        grup/.style={draw, dashed, gray!60, rounded corners, inner sep=5pt},
        panah/.style={-{Latex[length=2mm]}, thick},
    ]
        \foreach \i/\x/\t in {1/1.9/{ctr}, 2/4.9/{ctr$+1$}, 3/7.9/{ctr$+2$}} {
            \node[blok] (k\i) at (\x,0)    {\t};
            \node[enc]  (e\i) at (\x,-1.2) {$E_K$};
            \node[xor]  (x\i) at (\x,-2.4) {$\oplus$};
            \node[cip]  (c\i) at (\x,-3.7) {$C_{\i}$};
            \node[blok] (p\i) at (\x-2.0,-2.4) {$P_{\i}$};
            \draw[panah] (k\i) -- (e\i);
            \draw[panah] (e\i) -- (x\i);
            \draw[panah] (p\i) -- (x\i);
            \draw[panah] (x\i) -- (c\i);
            % Kotak putus-putus menandai satu blok yang diproses mandiri,
            % supaya jelas P_i hanya menuju XOR pada panelnya sendiri.
            \node[grup, fit=(k\i)(c\i)(p\i)] {};
        }
        \node[font=\scriptsize, align=center] at (4.9,-4.5)
            {tidak ada ketergantungan antar blok: seluruh blok dapat diproses secara paralel};
    \end{tikzpicture}
    \caption{Mode Counter (CTR): nilai counter yang terus bertambah dienkripsi lalu di-XOR dengan plainteks}
    \label{fig:mode-ctr}
\end{figure}
